\section{Carbon and Capacity}
\label{sec:carbon}

\subsection{Empty Miles as Emissions Without Purpose}

An empty truck emits approximately 85--90\% of the CO$_2$ per mile of a fully loaded truck. This counterintuitive result reflects the dominant contribution of engine load (not payload weight) to fuel consumption: a diesel engine running at 70 mph on a flat grade burns fuel primarily to overcome aerodynamic drag, which scales with vehicle frontal area---the same whether the trailer is full or empty. The EPA MOVES3 model reports a baseline emission factor of approximately 161.8 grams of CO$_2$ per mile for a Class 8 tractor-trailer at highway speed, dropping to approximately 148 grams per mile when fully loaded (improved aerodynamics from trailer weight lowering the center of gravity) \citep{EPA_MOVES3}.

For the purposes of the carbon accounting in this paper, we use a conservative 87\% factor: an empty truck emits 87\% of the CO$_2$ of a loaded truck for the same distance traveled. At 35\% empty miles and 300 billion total truck-miles nationally:

\begin{align}
  E_{\text{empty}} &= 0.35 \times 300 \times 10^9 \text{ miles} \times 148 \text{ g/mile} \times 0.87 \notag \\
  &= 105 \times 10^9 \text{ miles} \times 128.8 \text{ g/mile} \notag \\
  &= 13.5 \times 10^{12} \text{ g CO}_2 \notag \\
  &= 13.5 \text{ MMT CO}_2/\text{year}
\end{align}

A 15 percentage-point reduction to 20\% empty (a 43\% reduction in empty miles) eliminates:

\begin{equation}
  \Delta E = 45 \times 10^9 \text{ miles} \times 128.8 \text{ g/mile} = 5.8 \text{ MMT CO}_2/\text{year}
\end{equation}

For reference, the US transportation sector emits approximately 1,800 MMT CO$_2$ per year, of which medium and heavy-duty trucks account for approximately 440 MMT (24\%) \citep{EPA_GHG2023}. A 5.8 MMT reduction represents approximately 1.3\% of trucking sector emissions, or 0.3\% of total transportation emissions---modest in absolute terms but achieved at zero infrastructure cost beyond the relay hub software platform.

\subsection{Effective Capacity Equivalent}

The capacity implication of the empty-mile reduction is more significant than the carbon accounting. When 45 billion empty truck-miles are replaced by 45 billion loaded truck-miles, the effective freight capacity of the network increases by the same amount as if 45 billion new loaded truck-miles were added to the system. At the average long-haul payload of 19 tons \citep{ATRI_costs2024}:

\begin{equation}
  \Delta \text{Freight} = 45 \times 10^9 \text{ miles} \times \frac{19 \text{ tons}}{800 \text{ miles (avg haul)}} = 1.07 \text{ billion ton-trips}
\end{equation}

In ton-miles, the newly-loaded miles add:
\begin{equation}
  \Delta \text{Ton-miles} = 45 \times 10^9 \text{ miles} \times 19 \text{ tons} = 855 \text{ billion ton-miles}
\end{equation}

US total freight (all modes) is approximately 5,700 billion ton-miles annually \citep{BTS_freight2023}. An 855 billion ton-mile increase in truck capacity represents a 15\% increase in freight capacity---equivalent to building an entirely new Class 8 fleet operating on previously empty corridors. No new trucks are required; no new roads are needed. The efficiency gain comes entirely from information improvement.

\subsection{Daily Empty-Mile Reduction}

Converting the annual total to a daily figure for operational context: 45 billion fewer empty miles per year equals 123 million fewer empty truck-miles per day nationally. At an average empty-mile speed of 65 mph, this corresponds to approximately 1.9 million truck-hours per day of reclaimed productive capacity.

The 45M fewer empty miles per day figure cited in the abstract reflects the headline reduction from 35\% to 20\% on the full eligible truck population. The 123 million figure uses the same calculation basis; the discrepancy arises from rounding and the distinction between miles and the subset used in the abstract. We use 45M fewer empty miles per day as a conservative estimate applicable to hub-participating corridors only (50\% of total truck-miles), not the full long-haul fleet.

\subsection{No New Trucks Needed}

The freight capacity equivalent of the empty-mile reduction is a critical policy point: Congress has periodically debated increasing the national gross vehicle weight limit from 80,000 to 88,000 or 97,000 lbs to increase freight throughput per truck-trip. The weight increase would require bridge reinforcement investments estimated at \$48--\$120 billion nationally \citep{FHWA_Bridge2022}.

The relay hub empty-mile reduction achieves a larger capacity increase---an effective 855 billion ton-miles versus approximately 120 billion ton-miles from the weight increase---at zero infrastructure cost. The comparison is not precise (the weight increase improves loaded capacity; the empty-mile reduction recaptures wasted capacity), but it illustrates the magnitude: the information-architecture gain from relay hub pre-matching is a larger freight capacity investment than a decade of weight-limit policy debate has sought to unlock.

\subsection{Carbon Intensity Improvement}

The carbon intensity of the trucking sector, expressed as grams of CO$_2$ per revenue ton-mile, improves under relay hub pre-matching in two ways:

\begin{enumerate}
  \item \textbf{Fewer empty miles}: The same CO$_2$ is now allocated over more revenue ton-miles. At baseline (35\% empty), revenue ton-miles are generated on 65\% of total truck-miles. At 20\% empty, revenue ton-miles are generated on 80\% of total truck-miles---a 23\% increase in the denominator.

  \item \textbf{Modal efficiency improvement}: A fully loaded truck is more fuel-efficient per ton-mile than the same truck running at 70\% load (a common practical loading). Hub pre-matching enables better load consolidation, increasing average load factor toward the 85--90\% range observed at UPS and FedEx facilities.
\end{enumerate}

Combined, the carbon intensity improvement is approximately 12--15\% reduction in grams CO$_2$ per revenue ton-mile---the figure cited in claim C4.3. This metric is preferable to absolute emission reduction for policy comparison because it controls for freight volume growth: as the economy grows and freight demand increases, tons per mile will grow, but the intensity improvement persists regardless of volume.
